\begin{abstract}
Este trabalho apresenta a aplicação da meta-heurística \textit{Ant Colony System} (ACS) para resolver o problema de acasalamento otimizado em rebanhos bovinos, visando simultaneamente minimizar a coancestralidade entre os descendentes e maximizar o ganho genético esperado. O problema é modelado como um grafo bipartido que conecta machos e fêmeas, no qual cada aresta representa uma possível combinação de acasalamento ponderada por feromônios e heurísticas genéticas. A implementação proposta realiza pré-processamento dos dados, indexação eficiente da coancestralidade e construção de soluções probabilísticas respeitando restrições reprodutivas, como o limite máximo de uso dos machos. Experimentos conduzidos com diferentes limites de acasalamento demonstraram que valores mais altos tendem a aumentar o ganho genético, mas também elevam significativamente a coancestralidade. A normalização dos critérios e a aplicação de uma função de utilidade permitiram identificar o melhor equilíbrio entre os objetivos conflitantes, que neste estudo ocorreu com o limite de 15 acasalamentos por macho. Os resultados confirmam a adequação do ACS para decisões multicritério em manejo genético, oferecendo soluções eficientes e biologicamente plausíveis.
\end{abstract}
\vspace{1em}

\begin{abstract}
This work applies the \textit{Ant Colony System} (ACS) metaheuristic to solve the optimized mating problem in cattle herds, aiming simultaneously to minimize coancestry among offspring and to maximize expected genetic gain. The problem is modeled as a bipartite graph connecting males and females, where each edge represents a possible mating combination weighted by pheromones and genetic heuristics. The proposed implementation performs data preprocessing, efficient coancestry indexing, and probabilistic solution construction while respecting reproductive constraints such as the maximum usage limit of each male. Experiments conducted under different mating limits show that higher limits tend to increase genetic gain but also significantly raise coancestry. Normalization of the criteria and the use of a utility function enabled identifying the best balance between the conflicting objectives, which in this study occurred at a limit of 15 matings per male. The results confirm the suitability of ACS for multicriteria decision-making in genetic management, providing efficient and biologically meaningful solutions.
\end{abstract}

\begin{IEEEkeywords}
Ant Colony System, Otimização Combinatória, Coancestralidade, Melhoramento Genético, Problema de Acasalamento, Meta-heurísticas.
\end{IEEEkeywords}
